\documentclass[11pt]{article}
\usepackage[margin=1in]{geometry}
\usepackage{amsmath,amssymb,amsthm,mathtools,booktabs}
\usepackage[T1]{fontenc}
\usepackage{lmodern,microtype}
\usepackage{xurl}
\usepackage[colorlinks=true,linkcolor=blue,citecolor=blue,urlcolor=blue]{hyperref}
\newtheorem{theorem}{Theorem}[section]
\newtheorem{lemma}[theorem]{Lemma}
\newtheorem{proposition}[theorem]{Proposition}
\newtheorem{corollary}[theorem]{Corollary}
\theoremstyle{remark}
\newtheorem{remark}[theorem]{Remark}
\newcommand{\OO}{\mathcal O}
\newcommand{\FF}{\mathbb F}
\newcommand{\PP}{\mathbb P}
\newcommand{\EE}{\mathbb E}
\newcommand{\Mat}{\operatorname{Mat}}
\newcommand{\Alt}{\operatorname{Alt}}
\newcommand{\Sp}{\operatorname{Sp}}
\newcommand{\GL}{\operatorname{GL}}
\newcommand{\rank}{\operatorname{rank}}
\newcommand{\TV}{\operatorname{TV}}
\newcommand{\im}{\operatorname{im}}
\newcommand{\length}{\operatorname{length}}
\newcommand{\Emb}{\operatorname{Emb}}
\newcommand{\trans}{\mathsf T}
\newcommand{\qbin}[2]{\genfrac{[}{]}{0pt}{}{#1}{#2}_q}
\title{First-digit optimality for symplectic blocks\\over local rings}
\author{Research note for the ProveIt GRH transfer program}
\date{27 November 2026}
\begin{document}
\maketitle
\begin{abstract}
An observed block of a Haar symplectic matrix is generally distinguishable from a
matrix of independent uniform entries. We prove that, in a stable dimension
range, the optimal equal-prior distinction is made entirely from the first
residue digit, even when both source and target contain arbitrarily many
symplectic pairs. For source and target dimensions \(2b,2a\) in ambient
dimension \(2n\), it is enough that
\(n\ge a+b+\min(a,b)-1\). The optimal event is singularity of one residual
alternating form modulo the uniformizer, and the total variation distance is
exactly the same at every finite precision and over the complete ring. The
proof counts residue classes on which the Gram map is smooth and proves a
strict, precision-independent sign for the likelihood. A separate theorem gives
an exact Fourier formula and the sharp bounded-density threshold for arbitrary
alternating Gram matrices. All results include residue characteristic two.
They extend the single-pair local theorem in the supplied repository; they use
no Riemann-hypothesis assumption.
\end{abstract}

\section{Source, question, and contribution}

The source for this note is the symplectic portion of the article in
\path{Analysis/GeneralizedRiemannHypothesis}, pinned to ProveIt commit
\texttt{6d284fcf75bc3a91bb2e79ba7636f94724bff62c} \cite{ProveIt}.
Its section \emph{Spanning, frames, auxiliary characters and matrix statistics}
contains an exact finite-field compression law, a residual-Gram reduction,
and a complete local-ring theorem for a source consisting of one symplectic
pair. In its labels, these are \texttt{spn:thm:complete-law},
\texttt{spn:prop:residual}, and \texttt{spn:thm:padic}.
Question \texttt{spn:q:symplectic} explicitly asks for multi-pair local laws
and for frame types whose first-digit total variation persists at every
precision.
We also inspected the newer incoming packet
\emph{Optimal tests for nondegenerate symplectic compressions}
\cite{FiniteTests}. It proves the finite-field singular-residual test for
all nonoverlapping dimension offsets and gives a uniform limiting profile.
Thus neither the finite-field sign test nor its finite-field asymptotics
is claimed as new here. That packet explicitly leaves multi-pair finite
local rings as a further question.

Theorem~\ref{thm:main} answers that persistence question in a dimension range
allowing both numbers of pairs to grow. It is an exact equality, with no
asymptotic error and no dependence on the ramification index.
Theorem~\ref{thm:fourier} gives the accompanying exact local-density algorithm,
and Theorem~\ref{thm:regularity} determines precisely when these Gram laws
have a bounded density.
The argument is unconditional linear algebra and probability over local
rings. It uses the repository's transfer direction rather than a consequence
of its asserted Dirichlet GRH theorem.

The single-pair case below is already in the article. The new statement here
is simultaneous persistence for multi-pair source and target in the stated
range, with a proof by the combined contribution of kernel-zero and
kernel-one residue maps. Pirani's work \cite{Pirani} supplies related
Lie-algebra lifting methods for traces over local fields; the \(p\)-adic
random-matrix literature also treats corners of Haar \(\GL\) matrices
\cite{VanPeski}. Neither reference is used as a substitute for the proof
below. The source comparison and targeted literature search support a new
extension relative to the inspected repository; they do not certify global
publication priority.

\section{Statement and statistical meaning}

Let \(\OO\) be a complete discrete valuation ring, with uniformizer \(\pi\)
and residue field \(\FF_q\). Its fraction field may have characteristic zero
or positive characteristic, and \(q\) may be even. Put
\[
 R_e=\OO/\pi^e\OO\quad(e\ge1),\qquad R_\infty=\OO.
\]
All finite spaces carry uniform probability and all compact groups or additive
spaces carry normalized Haar probability. Let \(J_v\) be the block diagonal
alternating matrix with \(v\) blocks
\(\left(\begin{smallmatrix}0&1\\-1&0\end{smallmatrix}\right)\).
For \(0\le k\le v\), write
\[
 D(v,k)=\prod_{i=v-k+1}^{v}(1-q^{-2i}),\qquad D(v,0)=1.
\]
Choose symplectic frames \(Y\in\Mat_{2n,2a}(\OO)\) and
\(X\in\Mat_{2n,2b}(\OO)\):
\[
 Y^\trans J_nY=J_a,\qquad X^\trans J_nX=J_b.
\]
For uniform \(g\in\Sp_{2n}(R_e)\), let \(\mu_e\) be the law of
\(M=Y^\trans J_ngX\), and let \(\nu_e\) be the independent-entry law on
\(\Mat_{2a,2b}(R_e)\). The residual form is
\[
 H(M)=J_b+M^\trans J_a^{-1}M.
 \tag{2.1}\label{eq:residual}
\]
It is alternating, including in characteristic two. Define
\[
 E=\{M:\det\overline{H(M)}=0\},\qquad
 \beta_{q,a,b}=\nu_1(E^c),
 \tag{2.2}\label{eq:event}
\]
where a bar denotes reduction modulo \(\pi\).

\begin{theorem}[Exact first-digit optimality]\label{thm:main}
Let \(a,b\ge1\) and
\[
 n\ge a+b+\min(a,b)-1.
 \tag{2.3}\label{eq:range}
\]
Put
\[
 \alpha=\frac{D(n-a,b)}{D(n,b)}.
 \tag{2.4}\label{eq:alpha}
\]
Then, for every \(e=1,2,\ldots,\infty\),
\[
 \boxed{\quad
 \TV(\mu_e,\nu_e)=(1-\alpha)\beta_{q,a,b}
                =\TV(\mu_1,\nu_1).\quad}
 \tag{2.5}\label{eq:tv}
\]
The event \(E\) is optimal at all these precisions. At each finite precision
the likelihood is exactly \(\alpha<1\) on \(E^c\), and is strictly greater
than one on \(E\). Over \(\OO\), the two measures have a continuous
likelihood with the same strict signs.
\end{theorem}

Here \(\TV(\mu,\nu)=\sup_A|\mu(A)-\nu(A)|\).
Under equally likely alternatives, the largest possible success probability
of a binary decision rule is \((1+\TV(\mu,\nu))/2\). The theorem says
that the rule \emph{declare a symplectic block exactly when the residual
form is singular modulo \(\pi\)} already achieves this probability even
if every digit is available.

This conclusion is about optimal equal-prior testing, not full statistical
sufficiency. Higher digits need not be uniformly distributed, and they can
matter for unequal prior probabilities or unequal error costs.
The bound \eqref{eq:range} is sufficient; we do not claim it is the sharp
threshold for precision persistence.

\begin{remark}[Symmetry]
We prove the theorem when \(c=n-a\ge2b-1\).
Interchanging the frames replaces \(M\) by \(-M^\trans\), with the same law
after \(g\mapsto g^{-1}\). The alternative condition is \(n-b\ge2a-1\);
their union is \eqref{eq:range}.
If \(P_v=\prod_{i=1}^v(1-q^{-2i})\), then
\[
 \frac{D(n-a,b)}{D(n,b)}
  =\frac{P_{n-a}P_{n-b}}{P_{n-a-b}P_n}
  =\frac{D(n-b,a)}{D(n,a)}.
\]
The residual nonsingularity events also agree after transposition, by
\(\det(I+UV)=\det(I+VU)\). Thus the formula does not depend on which frame
is designated as the source.
\end{remark}

\section{Gram maps and smooth residue classes}

For this section put \(s=2b\) and \(m=\binom{s}{2}\).
For a uniform \(C\in\Mat_{2c,s}(R_e)\), define
\[
 G_{c,s,e}=C^\trans J_cC,\qquad
 f_{c,s,e}(H)=q^{me}\PP(G_{c,s,e}=H).
 \tag{3.1}\label{eq:gramdensity}
\]
Thus \(f\) is the density relative to uniform alternating matrices.

\begin{lemma}[The differential in every characteristic]\label{lem:diff}
For \(C_0\in\Mat_{2c,s}(\FF_q)\), put \(K=\ker C_0\) and
\(\ell=\dim K\). The linear map
\[
 L_{C_0}(Z)=C_0^\trans J_cZ+Z^\trans J_cC_0
 \tag{3.2}\label{eq:diff}
\]
has as image exactly the alternating forms vanishing on \(K\times K\).
Its image has codimension \(\binom{\ell}{2}\).
In particular it is surjective when \(\ell=0\) or \(1\).
\end{lemma}

\begin{proof}
Choose a source basis making \(C_0=[C'\;0]\), with \(C'\) injective.
The map \(C'^\trans J_c\) has full row rank, so its product with \(Z\)
can be prescribed freely. Consequently the upper-left block of
\eqref{eq:diff} is an arbitrary matrix minus its transpose, its cross-block
is arbitrary, and its bottom-right \(\ell\)-square block is zero.
Every alternating matrix is a matrix minus its transpose: choose one
triangle and set the other triangle and diagonal to zero.
This requires no division by two. It proves the asserted image and
codimension.
\end{proof}

\begin{corollary}[Flat lifting from a smooth residue class]\label{cor:lifts}
Suppose \(\dim\ker C_0\le1\) and \(C_0^\trans J_cC_0=H_0\).
For every lift \(H\) of \(H_0\) to \(R_e\), precisely
\[
 q^{(2cs-m)(e-1)}
 \tag{3.3}\label{eq:liftcount}
\]
lifts of \(C_0\) have Gram matrix \(H\).
Hence the residue class \(C_0\) contributes the constant \(q^{m-2cs}\)
to \(f_{c,s,e}(H)\), independent of every higher digit of \(H\).
\end{corollary}

\begin{proof}
At a lift from precision \(j\) to \(j+1\), replacing \(C\) by
\(C+\pi^jZ\) changes the new Gram digit by \(L_{C_0}(\overline Z)\).
The quadratic correction is divisible by \(\pi^{2j}\), hence by
\(\pi^{j+1}\). Lemma~\ref{lem:diff} gives exactly \(q^{2cs-m}\)
solutions for every required correction. Induction proves
\eqref{eq:liftcount}; multiplication by the normalization in
\eqref{eq:gramdensity} proves the final statement.
\end{proof}

\begin{lemma}[Embedding count]\label{lem:emb}
Fix an alternating form of rank \(2r\) on \(\FF_q^d\).
The number of injective maps taking this form to \(J_c\) is
\[
 \Emb_c(d,r)=q^{2cd-\binom d2}D(c,d-r)
 \quad\text{if }d-r\le c,
 \tag{3.4}\label{eq:emb}
\]
and is zero if \(d-r>c\).
\end{lemma}

\begin{proof}
Split the source into \(r\) symplectic pairs and a radical of dimension
\(t=d-2r\). Choosing the images of the pairs gives
\[
 \frac{|\Sp_{2c}(\FF_q)|}{|\Sp_{2(c-r)}(\FF_q)|}
   =q^{4cr-2r^2+r}D(c,r).
\]
The radical images are an ordered independent isotropic \(t\)-tuple
in a symplectic space of dimension \(2(c-r)\). After \(j\) choices,
the next vector lies in their orthogonal space but outside their span,
so its number of choices is \(q^{2(c-r)-j}-q^j\).
Multiplying these factors for \(0\le j<t\) proves both the obstruction
and \eqref{eq:emb}.
\end{proof}

\begin{lemma}[Likelihood reduction]\label{lem:likelihood}
For finite \(e\),
\[
 \frac{\mu_e(M)}{\nu_e(M)}
   =\frac{f_{c,s,e}(H(M))}{D(n,b)},\qquad c=n-a.
 \tag{3.5}\label{eq:likelihood}
\]
\end{lemma}

\begin{proof}
Symplectic Gram--Schmidt extends each prescribed frame to a symplectic
basis over \(R_e\), and the group acts transitively on frames of Gram
matrix \(J_b\). Write \(gX=(B;C)\) in the orthogonal decomposition
determined by \(Y\). Then \(M=J_aB\) and
\[
 C^\trans J_cC=J_b-B^\trans J_aB
              =J_b+M^\trans J_a^{-1}M.
\]
There is no extra injectivity condition: a matrix with Gram \(J_b\)
has a left inverse. By Lemma~\ref{lem:emb} and
Corollary~\ref{cor:lifts}, the total number of source frames is
\(q^{(2ns-m)e}D(n,b)\). For fixed \(B\), the number of admissible \(C\)
is \(q^{(2cs-m)e}f_{c,s,e}(H)\).
Dividing this ratio by \(\nu_e(M)=q^{-2ase}\) proves the claim.
\end{proof}

\section{A sign that cannot change at higher precision}

\begin{proposition}[Uniform strict lower bound]\label{prop:sign}
Assume \(s=2b\ge2\) and \(c\ge s-1\).
If the reduction of \(H\in\Alt_s(R_e)\) is singular, then
\[
 f_{c,s,e}(H)>1
 \tag{4.1}\label{eq:strictsign}
\]
for every finite \(e\). The lower bound is independent of \(e\).
If \(\overline H\) is nonsingular, then
\[
 f_{c,s,e}(H)=D(c,b).
 \tag{4.2}\label{eq:nonsingsign}
\]
\end{proposition}

\begin{proof}
Let \(\rank\overline H=2r\) and \(t=s-2r\). For a singular form,
\(t\ge2\). The kernel of every matrix with Gram \(\overline H\)
is contained in its radical. There are
\([t]_q=(q^t-1)/(q-1)\) choices of a kernel line.
The contribution of matrices with kernel zero or one, using the preceding
lemmas, gives
\[
 f_{c,s,e}(H)\ge
 q^{m-2cs}\bigl(\Emb_c(s,r)+[t]_q\Emb_c(s-1,r)\bigr).
 \tag{4.3}\label{eq:smoothlower}
\]
Put \(u=c-s+r+1\). Our assumption gives \(u\ge0\).
Substitution into \eqref{eq:smoothlower} simplifies its right side to
\[
 L=
 \prod_{j=u+1}^{c}(1-q^{-2j})
 \left(1+q^{-2u}\sum_{j=1}^{t-1}q^{-j}\right).
 \tag{4.4}\label{eq:L}
\]
For \(u=0\), the zero-kernel contribution is zero and the kernel-one
term alone gives the same expression.

To show \(L>1\), set \(A=q^{-2u}\le1\). The finite-product inequality gives
\[
 \prod_{j=u+1}^{c}(1-q^{-2j})
 \ge1-\sum_{j=u+1}^{c}q^{-2j}
 >1-\frac{A}{q^2-1}.
\]
Since \(t\ge2\),
\[
 L>
 \left(1-\frac{A}{q^2-1}\right)\left(1+\frac A q\right)
 =1+\frac{A(q^2-q-1-A)}{q(q^2-1)}\ge1.
 \tag{4.5}\label{eq:Lpositive}
\]
The final inequality uses \(q\ge2\) and \(A\le1\); equality in its numerator
can occur at \(q=2,A=1\), but the preceding strict inequality still yields
\(L>1\). This proves \eqref{eq:strictsign} also at the stable boundary.

If \(\overline H\) is nonsingular, every matrix with that Gram form is
injective. There are \(q^{2cs-m}D(c,b)\) such residue matrices and all
their higher lifts are flat. There are no other residue classes to count,
which proves \eqref{eq:nonsingsign}.
\end{proof}

\begin{proof}[Proof of Theorem~\ref{thm:main}]
By symmetry it suffices to assume \(c=n-a\ge2b-1\).
Since \(a>0\), one has \(D(c,b)<D(n,b)<1\).
The likelihood formula and Proposition~\ref{prop:sign} show that the
likelihood is \(\alpha<1\) on \(E^c\), and is strictly greater than one
on \(E\), uniformly in precision. The negative part of
\(\mu_e-\nu_e\) is therefore exactly
\((1-\alpha)\nu_e|_{E^c}\). Its mass is
\((1-\alpha)\beta_{q,a,b}\), proving the finite-precision conclusion.

The inverse-limit matrix space is compact and its cylinder sets generate
the Borel sigma-algebra. Total variation is the supremum of total
variations on its finite quotients: for instance, approximate a Borel
indicator by its conditional expectations with respect to the finite
quotients in \(L^1(\mu_\infty+\nu_\infty)\). Consequently the same
distance holds for \(e=\infty\), and the fixed cylinder event \(E\)
attains it. Alternatively, the measure inequality on \(E\) and the
exact restricted-measure identity on \(E^c\) pass from cylinders by a
monotone-class argument. The continuous likelihood and strict signs
follow from Theorem~\ref{thm:regularity} and the
precision-independent lower bound \eqref{eq:L}.
\end{proof}

\section{Exact local densities and their regularity threshold}

This section concerns arbitrary \(s\ge2\), not necessarily even.
It gives an exact way to compute the densities without enumerating all
\(2c\)-by-\(s\) input matrices.

\begin{theorem}[Finite Fourier formula]\label{thm:fourier}
Let \(\psi_e\) be a primitive additive character of \(R_e\), and put
\(\langle W,H\rangle=\sum_{i<j}W_{ij}H_{ij}\).
For every \(H\in\Alt_s(R_e)\),
\[
 f_{c,s,e}(H)=
 \sum_{W\in\Alt_s(R_e)}
 q^{-c\length_\OO(\im W)}
 \psi_e(-\langle W,H\rangle).
 \tag{5.1}\label{eq:fourier}
\]
The right side is an exact finite sum; it is not asserted to be a
closed product in the elementary divisors of \(H\).
\end{theorem}

\begin{proof}
Group the rows of \(C\) into \(c\) independent pairs \(x,y\).
The Gram contribution of a pair is \(xy^\trans-yx^\trans\), whose
pairing with \(W\) is \(x^\trans Wy\).
For a fixed \(x\), averaging \(\psi_e(x^\trans Wy)\) over \(y\)
gives one exactly when \(Wx=0\), and zero otherwise.
The perfect additive pairing on \(R_e\) is valid because \(\psi_e\)
is primitive. Averaging over \(x\) gives
\(|\ker W|/|R_e|^s=q^{-\length_\OO(\im W)}\).
Independence and finite Fourier inversion prove \eqref{eq:fourier}.
\end{proof}

Define
\[
 Z_e(c,s)=\sum_{W\in\Alt_s(R_e)}q^{-c\length_\OO(\im W)},
 \qquad Z_0(c,s)=1.
\]
For \(s=0,1\) set every \(Z_e(c,s)=1\).
Let \(A_s(r)\) be the number of alternating \(s\)-square matrices
of rank \(2r\) over \(\FF_q\):
\[
 A_s(r)=\qbin{s}{2r}\frac{|\GL_{2r}(\FF_q)|}{|\Sp_{2r}(\FF_q)|}.
 \tag{5.2}\label{eq:Ar}
\]
This is obtained by choosing the radical and then a nonsingular
alternating form on the quotient.

\begin{lemma}[Conductor recursion]\label{lem:recursion}
For \(e\ge1\),
\begin{align}
 Z_e(c,s)-Z_{e-1}(c,s)
  =\sum_{r=1}^{\lfloor s/2\rfloor}
  &A_s(r)q^{-2cr}
    q^{-r(2c-2s+2r+1)(e-1)}
    Z_{e-1}(c,s-2r).
 \tag{5.3}\label{eq:recursion}
\end{align}
\end{lemma}

\begin{proof}
The matrices divisible by \(\pi\) contribute \(Z_{e-1}(c,s)\).
For a fixed residue matrix of rank \(2r>0\), a fixed integral
change of basis makes its upper-left \(2r\)-square block invertible.
Writing the other blocks as \(B,D\), congruence eliminates the cross-block
and leaves the alternating Schur complement
\(D+B^\trans A^{-1}B=\pi T\).
For each \(T\in\Alt_{s-2r}(R_{e-1})\), the choices of \(A,B\)
number
\[
 q^{(e-1)(\binom{s}{2}-\binom{s-2r}{2})}.
\]
The image length is \(2re+\length_\OO(\im T)\).
Multiplying, summing over \(T\), and then over the \(A_s(r)\) residue
matrices yields \eqref{eq:recursion}.
Alternation of the Schur complement holds in characteristic two as well:
the inverse of a nonsingular alternating matrix defines an alternating
form.
\end{proof}

\begin{theorem}[Sharp bounded-density threshold]\label{thm:regularity}
Let \(s\ge2\) and \(c\ge0\) be integers.
For \(C\) additive-Haar in \(\Mat_{2c,s}(\OO)\), the law of
\(C^\trans J_cC\) has a bounded Haar density if and only if \(c\ge s-1\).
In that range it has a continuous density \(f_{c,s,\infty}\), obtained
by uniform convergence of the finite densities lifted to \(\Alt_s(\OO)\).
Writing \(Z_\infty(c,s)=\lim_eZ_e(c,s)\), one has the exact error
\[
 \|f_{c,s,\infty}-f_{c,s,e}\|_\infty
       =Z_\infty(c,s)-Z_e(c,s).
 \tag{5.4}\label{eq:uniformerror}
\]
For fixed \(q,c,s\) in this range, let \(\lambda=2c-2s+3>0\).
Then
\[
 Z_\infty(c,s)-Z_e(c,s)
 \sim
 \frac{A_s(1)q^{-2c}Z_\infty(c,s-2)}
      {1-q^{-\lambda}}\,q^{-\lambda e}.
 \tag{5.5}\label{eq:sharpdecay}
\]
\end{theorem}

\begin{proof}
Choose compatible primitive characters at all levels. In the Fourier
series on \(\Alt_s(\OO)\), the contribution of characters of exact
conductor \(e\) has total absolute coefficient mass
\(Z_e(c,s)-Z_{e-1}(c,s)\).
If \(c\ge s-1\), the exponent
\[
 r(2c-2s+2r+1)\ge r(2r-1)
\]
in \eqref{eq:recursion} is positive.
Induction on \(s\), using the recursion and geometric summation, shows
that the \(Z_e(c,s)\) are bounded. The Fourier series therefore converges
absolutely and uniformly to a continuous function. Its finite-quotient
integrals agree with those of the Gram law, proving that it is its density.
The triangle inequality bounds the omitted tail by
\(Z_\infty-Z_e\); at \(H=0\) all characters equal one, so equality
holds, proving \eqref{eq:uniformerror}.

The \(r=1\) summand of \eqref{eq:recursion} is
\[
 A_s(1)q^{-2c}q^{-\lambda(e-1)}Z_{e-1}(c,s-2).
\]
Its final factor converges to a positive finite limit. Every \(r\ge2\)
summand has a strictly larger exponential decay rate. Summing the resulting
geometric tail proves \eqref{eq:sharpdecay}.

Conversely, if \(c\le s-2\), the \(r=1\) summand alone gives
\[
 Z_e(c,s)-Z_{e-1}(c,s)
 \ge A_s(1)q^{-2c}q^{-(2c-2s+3)(e-1)},
\]
which tends to infinity. Thus \(f_{c,s,e}(0)=Z_e(c,s)\) is unbounded.
These values are the average density on the shrinking neighborhoods
\(\pi^e\Alt_s(\OO)\). An essentially bounded density would bound all
these averages, a contradiction.
\end{proof}

\begin{remark}
The density threshold \(c=s-1\) is sharp.
This does not establish that the testing threshold in
Theorem~\ref{thm:main} is sharp. A measure with an unbounded density can
still have its likelihood above one throughout the relevant residue
classes. The two questions are different.
\end{remark}

\section{Computing the total variation constant}

The number \(\beta_{q,a,b}\) in the main theorem involves only the
residue field and is completely explicit.
Set \(s=2b\) and define
\[
 \mathcal Z_a(s,r)=
 \sum_{\ell=0}^{s-2r}
 \qbin{s-2r}{\ell}\Emb_a(s-\ell,r).
 \tag{6.1}\label{eq:Zcount}
\]
Terms that violate the embedding range are zero. The formula chooses
the kernel inside the radical of a prescribed rank-\(2r\) form and
then embeds the quotient. Thus it counts all matrices whose Gram
matrix is that prescribed form. Consequently,
\[
 \boxed{\quad
 \beta_{q,a,b}
 =q^{-2as}
 \sum_{\substack{G\in\Alt_s(\FF_q)\\\det(J_b-G)\ne0}}
      \mathcal Z_a\bigl(s,\rank(G)/2\bigr).\quad}
 \tag{6.2}\label{eq:betaexact}
\]
This formula enumerates \(q^{\binom s2}\) alternating matrices rather
than \(q^{2as}\) observed matrices, and is independent of precision.
Together with \eqref{eq:alpha}, it is an exact finite formula for the
answer.

There is an even cheaper computation inherited from \cite{FiniteTests}.
Start a Markov chain at state \(b\) on \(\{0,\ldots,b\}\).
At state \(i\), put \(d=s-2i\) and take the transition probabilities
\[
 p_i^+=(1-q^{-d})(1-q^{1-d}),\qquad
 p_i^-=q^{-2d-1}(1-q^{-2i}),\qquad
 p_i^0=1-p_i^+-p_i^- .
 \tag{6.3}\label{eq:chain}
\]
The factors make \(p_b^+=p_0^-=0\).
After \(a\) steps, the mass at \(b\) is precisely
\(\beta_{q,a,b}\). Thus the constant in the local-ring theorem can be
computed in \(O(ab)\) rational arithmetic operations and \(O(b)\) storage,
independently of \(e\).

For completeness, adding one independent pair of rows to a Gram matrix
adds a uniform decomposable form \(uv^\trans-vu^\trans\).
At half-rank \(i\), an increase occurs when the images of \(u,v\) in the
\(d\)-dimensional quotient by the current column space are independent,
giving \(p_i^+\). A decrease requires both vectors in the column space,
with probability \(q^{-2d}\); the first must be nonzero there, and then
one prescribed nonzero pairing with the second cancels a symplectic
plane, giving \(q^{-1}(1-q^{-2i})\).
These are \(p_i^-\). A rank-two update changes half-rank by at most one,
so the remaining probability is \(p_i^0\). Starting at the nonsingular
form \(J_b\) and adding \(a\) pairs proves the claim.

\subsection{The first example with two pairs on both sides}

Take \(q=2,\ a=b=2,\ c=3,\ n=5\). Then an observed block has size
\(4\) by \(4\), inside a symplectic \(10\)-dimensional space, and
\[
 \alpha=\frac{5376}{5797},\qquad
 \beta_{2,2,2}=\frac{469}{1024}.
\]
Therefore
\[
 \TV(\mu_e,\nu_e)
 =\frac{197449}{5936128}
 =0.033262254452734175\ldots
 \tag{6.4}\label{eq:example}
\]
for \(e=1,2,1000,\infty\), and all other positive precisions.
The residue probability \(\beta\) was also verified independently by
enumerating all \(65\,536\) binary observed matrices: exactly
\(30\,016\) have a nonsingular residual.
The optimal equal-prior success probability is approximately
\(0.5166311272263671\), regardless of the number of available digits.

\begin{center}
\small
\begin{tabular}{ccccc}
\toprule
\(q\)&\((a,b,c)\)&\(\alpha\)&\(\beta_{q,a,b}\)&all-precision \(\TV\)\\
\midrule
\(2\)&\((2,2,3)\)&\(5376/5797\)&\(469/1024\)&\(197449/5936128\)\\
\(3\)&\((2,2,3)\)&\(597051/605242\)&\(38089/59049\)&\(7609439/871681338\)\\
\(3\)&\((1,2,3)\)&\(81/82\)&\(163/243\)&\(163/19926\)\\
\bottomrule
\end{tabular}
\end{center}

\section{Verification, scope, and reproducibility}

The proof is human-readable and complete; it is not a Lean formalization.
The included program \texttt{code/symplectic\_precision\_verify.py}
uses standard-library Python, exact integers, and rational arithmetic.
Floating-point arithmetic is used only to display decimal approximations.
Run
\begin{verbatim}
python3 code/symplectic_precision_verify.py
\end{verbatim}
from the archive directory. The default output is
\texttt{data/symplectic\_precision\_checks.json}.

The checks have independent components:
\begin{itemize}
\item Integer Fourier transforms compute all Gram likelihoods for the
  \(4\)-column examples over \(\mathbb Z/2\), \(\mathbb Z/3\),
  \(\mathbb Z/4\), and \(\FF_2[t]/(t^2)\).
  The last two rings distinguish mixed and equal characteristic at
  precision two.
\item Direct enumeration of one-pair Gram inputs independently checks the
  Fourier computation, with \(137\,889\) ordered pairs in total.
\item Every one of \(9\,714\) residual alternating matrices in the displayed
  cases is checked for the predicted likelihood sign, normalization, and
  exact total variation. In particular, all \(4\,096\) precision-two
  residuals are checked separately in each of the two characteristic-two
  local rings.
\item The binary \(4\)-by-\(4\) observed matrices and the ternary
  \(2\)-by-\(4\) observed matrices are independently enumerated to check
  \(\beta\), with \(65\,536+6\,561=72\,097\) matrices in total.
\item The strict lower bound is checked in \(1\,728\) exact rational
  parameter cases, and the conductor recursion is checked against direct
  image-length enumeration in \(60\) cases.
\end{itemize}
These checks test normalizations, characteristic-two conventions, and the
finite formulas. They are not the proof of the infinite family.

The remaining mathematical scope is clear. The theorem covers every
complete DVR in the stable range; the proof does not settle the optimal
testing range below it. The finite Fourier formula is exact at arbitrary
precision but its cost grows with the number of alternating coordinates.
A closed expression in the residual form's elementary divisors could
give a more efficient algorithm. Nor do we identify the likelihood-ratio
tests for every unequal-prior decision problem. Those questions remain
available for further work.

\begingroup
\small
\raggedright
\begin{thebibliography}{9}
\setlength{\itemsep}{2pt}
\bibitem{ProveIt}
Vladimir Reshetnikov et al.,
\emph{Beyond the Riemann hypothesis for Dirichlet \(L\)-functions:
consequences, transfers and limits},
ProveIt, section \emph{Spanning, frames, auxiliary characters and matrix statistics},
commit \texttt{6d284fcf75bc3a91bb2e79ba7636f94724bff62c}.
\url{https://github.com/VladimirReshetnikov/ProveIt/blob/6d284fcf75bc3a91bb2e79ba7636f94724bff62c/Analysis/GeneralizedRiemannHypothesis/Article/sections/spn.tex}.
\bibitem{Pirani}
Noam Pirani,
\emph{Traces of powers of random matrices over local fields},
arXiv:2408.04061, 2024.
\url{https://arxiv.org/abs/2408.04061}.
\bibitem{FiniteTests}
\emph{Optimal tests for nondegenerate symplectic compressions},
research packet \texttt{grh-two-open-questions.zip},
member \texttt{symplectic/symplectic.tex}, 27 November 2026.
ProveIt incoming archive at the same pinned commit.
\url{https://github.com/VladimirReshetnikov/ProveIt/blob/6d284fcf75bc3a91bb2e79ba7636f94724bff62c/Analysis/GeneralizedRiemannHypothesis/docs/incoming/grh-two-open-questions.zip}.
\bibitem{VanPeski}
Roger Van Peski,
\emph{Local limits in \(p\)-adic random matrix theory},
Proceedings of the London Mathematical Society.
\url{https://doi.org/10.1112/plms.12626}.
\end{thebibliography}
\endgroup
\end{document}
